\chapter{The Main page}
\label{ch:main}

The Main page tests a device in three steps, from top to bottom:
\emph{Connection} (how to reach it), \emph{Request} (what to ask) and
\emph{Test} (send, and see the answers).

\screenshot[8cm]{main-page}{The Main page: Connection, Request and Test.}
  {The Main page with a manual request: step 1 set to Modbus TCP, step 2 reading
   10 holding registers, step 3 after a few cycles with the Read registers tab.}

\section{Step 1: Connection}

\subsection{Interface}

Choose how the device is reached:

\begin{center}
\begin{tabularx}{\linewidth}{l L}
  \toprule
  \thead{Interface} & \thead{Use it for} \\
  \midrule
  Serial, RTU & RS-485 or RS-232 devices: the usual Modbus on a serial line. \\
  Serial, ASCII & Older devices that send frames as text, checked with an LRC. \\
  Modbus TCP & Devices and gateways on Ethernet (port 502). \\
  Modbus UDP & Devices that answer Modbus TCP frames over UDP. \\
  RTU over TCP & Serial-to-Ethernet converters that pass RTU frames as they are. \\
  RTU over UDP & The same, over UDP. \\
  \bottomrule
\end{tabularx}
\end{center}

\subsubsection{Serial}

\field{Serial port} lists the ports of this computer (\ui{↻} looks again
after plugging an adapter in). Set \field{Baud rate}, \field{Data bits},
\field{Parity} and \field{Stop bits} as the device is configured: every device
on the line must use the same.

\subsubsection{Network}

\field{IP address} and \field{Port}. The port is 502 unless the device or the
converter uses another.

\subsection{Device}

\field{Device (unit ID)} is the device's address: \val{1}--\val{247} on a serial
line or behind a gateway. \field{Timeout (ms)} is how long to wait for an
answer before counting a timeout; \field{Delay (ms)} is the pause between two
requests, which slow devices and some converters need.

\begin{tip}[Not sure of the settings?]
\ui{Find…} scans for unit IDs that answer, or, on a serial line, for the baud
rate, parity and stop bits of a device (Section~\ref{sec:find-devices}). The
Statistics tool suggests a timeout from the measured response times
(Section~\ref{sec:statistics}).
\end{tip}

\section{Step 2: Request}

\field{Source} chooses between a \emph{manual request} and a \emph{device map}.

\subsection{Manual request}

\begin{description}[style=nextline, leftmargin=0pt, font=\intermedium]
  \item[Command] The function code: read coils (FC01), discrete inputs (FC02),
    holding registers (FC03) or input registers (FC04), or one of the writes
    (FC05, FC06, FC15, FC16).
  \item[Start address] The first address, in the notation chosen in Settings
    (below).
  \item[Quantity] How many coils or registers. Modbus allows up to 125
    registers or 2000 bits per read.
\end{description}

\subsection{Device map}

With \field{Source} set to a map, \field{Map} chooses one of the maps of the
library, an open file, or \nav{From a file…}. Modbench then reads every
register of the map by name (Chapter~\ref{ch:maps}). The map's connection
defaults (unit ID, baud rate, parity, port) fill step 1 when the map is
picked; you can still change them.

\subsection{Byte order}

\field{Byte order} reads every value of a map with another byte and word order
than the map says. It is meant for checking: if the values only make sense
with another order, the map, or the device's configuration, is wrong.

\subsection{Loop}
\label{sec:loop}

With \field{Loop command} on, the request repeats until \ui{Stop}. The time
next to it is from the start of one cycle to the start of the next; \val{0}
repeats as soon as \field{Delay} allows. It can be changed while the test runs.
Without Loop, \uiaccent{Start} sends the request once.

\section{Address notation}
\label{sec:notation}

Manuals write addresses in many ways. The address menu on the
\tab{Device map}, \tab{Read registers} and \tab{Log} tabs chooses how Modbench
shows them:

\begin{center}
\begin{tabularx}{\linewidth}{l l L}
  \toprule
  \thead{Notation} & \thead{Holding register 100} & \thead{Notes} \\
  \midrule
  0-based & \val{100} & The address on the wire. \\
  1-based & \val{101} & As many manuals number registers. \\
  Hex & \val{0x0064} & As in protocol documents. \\
  Modicon & \val{400101} & Area digit in front (\val{0}, \val{1}, \val{3}, \val{4}), 1-based. \\
  IEC 61131-3 & \val{\%MW100} & As PLC programs write them. \\
  \bottomrule
\end{tabularx}
\end{center}

Address fields also take the Modicon and IEC forms as typed, and switch the
command to their area: typing \val{300010} reads input register 9.

\section{Step 3: Test}

\uiaccent{Start} sends the request (once, or in a loop); \ui{Stop} ends it. The
counters show \field{Valid responses}, \field{Error responses} (exceptions)
and \field{Timeouts}; the button next to each counter sets it back to zero.

\screenshot[2.5cm]{test-controls}{Start, Stop and the counters.}
  {The top of step 3 during a test: Start, Stop and the three counters with
   some values.}

The tabs below show the results:

\begin{description}[style=nextline, leftmargin=0pt, font=\intermedium]
  \item[\tab{Device map} or \tab{Read registers}] The values: by register name
    with a map (Chapter~\ref{ch:maps}), as a grid for a manual read
    (Chapter~\ref{ch:read-write}).
  \item[\tab{Chart}] Values over time (Chapter~\ref{ch:charts}).
  \item[\tab{Write values}] Write coils and registers
    (Section~\ref{sec:write}).
  \item[\tab{Raw data}] The bytes of each frame, decoded
    (Section~\ref{sec:raw}).
  \item[\tab{Log}] Every request with its result and time
    (Section~\ref{sec:log}).
\end{description}

\subsection{When the link drops}

With \nav{Reconnect when the connection is lost} on (in Settings), a test that
loses its link (a cable pulled, a gateway restarted) keeps trying to
reconnect, and goes on where it stopped.

\section{Finding devices}
\label{sec:find-devices}

\ui{Find…} in step 1 opens a scan:

\begin{description}[style=nextline, leftmargin=0pt, font=\intermedium]
  \item[\tab{Unit IDs}] Asks every unit ID in a range (\val{1}--\val{247} by
    default) and lists the ones that answer. Use a short \field{Wait (ms)} to
    go faster on a quiet line.
  \item[\tab{Serial settings}] For a serial device whose settings are
    unknown: tries the chosen baud rates, parities and stop bits, most common
    first, until the device answers.
\end{description}

\screenshot[5cm]{find-devices}{Finding unit IDs and serial settings.}
  {The Find panel after a unit ID scan that found two devices, with the
   Serial settings tab visible.}

\ui{Use} next to a unit ID that answered puts it in step 1.
